\documentclass[12pt]{article}
\usepackage[utf8]{inputenc}
\usepackage[T1]{fontenc}
\usepackage{hyperref}
\title{Sanggau Language Resources for K-12 Curriculum Development}
\author{Compiled by Claude Sonnet 5.5 (claude-sonnet-5-5) for Andr\'as Kornai}
\date{}
\begin{document}
\maketitle

\section{Introduction}
Sanggau (ISO 639-3 scg; Glottocode sang1339) is a Land Dayak language of Sanggau Regency, West Kalimantan, Indonesia. Wikipedia gives about 45,000 speakers, from a 1981 count, and says its varieties are quite divergent and may be distinct languages. ASJP gives about 31,300 speakers. It must not be confused with Kembayan (Matéq), a neighbouring Land Dayak language with its own sketch grammar. Very little exists for Sanggau, and no school materials were found.

\section{Primary Educational Resources}
None found. No local-curriculum (muatan lokal) materials for Sanggau turned up.

\section{Linguistic and Reference Materials}
\begin{itemize}
\item Sanggau language (Wikipedia): \url{https://en.wikipedia.org/wiki/Sanggau_language}
\item Land Dayak languages (Wikipedia): \url{https://en.wikipedia.org/wiki/Land_Dayak_languages}
\item ASJP Sanggau word list: \url{https://asjp.clld.org/languages/SANGGAU}
\item Numerals data, Sanggau (MPI-EVA, from Smith 2017): \url{https://lingweb.eva.mpg.de/channumerals/Sanggau.htm}
\end{itemize}
No dictionary or grammar of Sanggau was found.

\section{Religious and Community Resources}
\begin{itemize}
\item Joshua Project, Sanggau-Sekadau people group: \url{https://joshuaproject.net/people_groups/14723}
\end{itemize}

\section{Digital and Audio Resources}
\begin{itemize}
\item Global Recordings Network, Sanggau audio-visual Bible lessons (evangelistic, not educational): \url{https://globalrecordings.net/en/language/scg}
\end{itemize}

\section{Curriculum Development Context}
Sanggau has no known writing standard and is not known to be taught in schools. The language is spoken in every generation, which makes documentation and curriculum work feasible. Indonesian is the language of instruction. A curriculum would need dialect decisions, since the varieties differ, together with an orthography, a lexicon and recorded texts from the community. Indonesian local-curriculum rules could provide a route into schools.

\end{document}